\subsection{A System of Three Equations by Substitution} 
Occasionally, we encounter a system that is set up to make substitution especially easy. These arise in some common applications.
\begin{Example}
Alice, Bob, and Carol are splitting a dinner bill. Bob ordered an expensive meal, so he will pay \$5 more than Alice. Carol is paying for Dave's meal as well, so she will pay twice as much as Alice. The total bill is \$65. How much will each person pay?
\begin{itemize}
\item Start by defining some variables:
\begin{variables}
A&the amount Alice pays\\
B&the amount Bob pays\\
C&the amount Carol pays
\end{variables}
\item Our equations are:
\[\begin{systemL}
\makeblank&\\
\makeblank&\\
\makeblank&
\end{systemL}\]
\item This is set up to make substitution easy, because the first two equations express \makeblanksmall and \makeblanksmall in terms of \makeblanksmall.
\end{itemize}
\begin{lpic}{}{10}
\foreach \y/\name in {-8/Alice,-9/Bob,-10/Carol} \regtextleft{1}{\y}{\name\ pays:};
\end{lpic}
\end{Example}
